% ===================================================================== Chapter 2
\chapter{Engineering System}\label{ch:system}

\section{System Description}
The system is a straight, thick-walled circular duct of Inconel 718 with air flowing through its bore. A uniform heat flux enters the outer
surface, conducts through the wall and is removed by the turbulent air flow. The end faces are adiabatic. Structurally, the duct either
expands freely (LC1) or is held axially at both ends (LC2), in which case the prevented thermal growth produces an axial compressive
force. All dimensions and loads are selected class-B inputs (Chapter~\ref{ch:params}).

\section{Heated Duct Geometry}
Figure~\ref{fig:drawing} shows the dimensioned drawing prepared in Section 3 of the project. The bore diameter of 20~mm and the length of
600~mm give $L/D_i$ = 30, long enough for the flow to become hydrodynamically developed over the rear part of the duct. The diameter ratio
$D_o/D_i$ = 2 produces a 10~mm wall, thick enough for a measurable through-wall temperature difference and for a stiff column section.

\begin{figure}[htbp]
\centering
\includegraphics[width=0.74\linewidth]{ENGINEERING_DRAWING.png}
\caption[Engineering drawing of the heated duct]{Engineering drawing of the heated duct: longitudinal section and cross-section (dimensions in mm).}
\label{fig:drawing}
\src{Source: \texttt{03\_CAD\_Geometry/Drawings/ENGINEERING\_DRAWING.png} (project drawing of the analysed duct).}
\end{figure}

\section{Geometry Parameters}
Table~\ref{tab:geometry} lists the geometry. The heated area differs between the true cylinder and the CFD mesh, because the CFD section is
an inscribed 48-sided polygon; the resulting 0.071\,\% difference in heat input is exact and is discussed in Chapters~\ref{ch:cfdres}
and~\ref{ch:mi}.

%% generated by gen_tables.py - RE-ANALYSIS 2026
\begin{table}[htbp]
\centering
\small
\caption{Duct geometry of the baseline (selected class-B inputs, verified in the CAD model)}
\label{tab:geometry}
\begin{tabular}{>{\raggedright\arraybackslash}p{0.42\linewidth}lrl}
\toprule
Parameter & Symbol & Value & Unit \\
\midrule
Inner (bore) diameter & $D_i$ & 20 & mm \\
Outer diameter & $D_o$ & 40 & mm \\
Wall thickness & $t$ & 10 & mm \\
Length & $L$ & 600 & mm \\
Hydraulic diameter & $D_h$ & 20 & mm \\
Slenderness of the flow passage & $L/D_i$ & 30 & -- \\
Diameter ratio & $D_o/D_i$ & 2 & -- \\
Fluid flow area & $A_c$ & 314.159 & mm$^2$ \\
Solid volume & $V_s$ & $5.654867\times10^{-4}$ & m$^3$ \\
Heated outer area, true cylinder & $A_o$ & 0.07539822 & m$^2$ \\
Heated outer area, CFD 48-facet polygon & $A_{o,48}$ & 0.07534441 & m$^2$ \\
\bottomrule
\end{tabular}
\par\smallskip\parbox{\linewidth}{\footnotesize\raggedright Source: \texttt{MASTER\_PROJECT\_DATA.csv} rows M023--M031 (CAD values from \texttt{cad\_measurements.json}; the 48-facet area from the Fluent surface-integral report). Status VERIFIED.}
\end{table}


\section{Fluid Domain}
The fluid domain is a solid body filling the bore ($\diameter$20~mm $\times$ 600~mm, volume $1.884956\times10^{-4}$ m$^3$). Air enters
through \texttt{FLUID\_INLET} at $z = 0$ and leaves through \texttt{FLUID\_OUTLET} at $z$ = 600~mm; its only lateral boundary is the
fluid--solid interface.

\section{Solid Domain}
The solid domain is the Inconel 718 annulus between $r$ = 10~mm and $r$ = 20~mm (volume $5.654867\times10^{-4}$ m$^3$, mass 4.631 kg). Its outer
cylinder carries the heat flux; its inner cylinder is the conjugate interface; its end annuli are adiabatic and, in the structural model,
carry the axial supports.

\section{Heat-Transfer Path}
Heat follows a single series path (Figure~\ref{fig:physics}): conduction radially inward through the wall, convection from the bore
surface into the air, and advection of the absorbed energy to the outlet. The analytical baseline shows that the conduction resistance of
the wall is only 3.4\,\% of the total resistance between outer wall and air; the air film carries the remaining 96.6\,\%
($R_{cond}/R_{conv}$ = 0.0342). The metal temperature is therefore controlled mainly by the heat-transfer coefficient on the air side, not by
the conductivity of the wall.

\begin{figure}[htbp]
\centering
\includegraphics[width=0.80\linewidth]{PHYSICS_SCHEMATIC.png}
\caption[Physics chain mapped onto the CAD entities]{Physics chain mapped onto the CAD entities: flow, convection, conduction, expansion and stress. The values printed in the five
process boxes are Section 2 analytical estimates, not final results.}
\label{fig:physics}
\src{Source: \texttt{03\_CAD\_Geometry/Drawings/PHYSICS\_SCHEMATIC.png} (Section 3).}
\end{figure}

\section{Structural Load Path}
The wall temperature rises from the inlet to the outlet by more than 100~K and exceeds the 300~K stress-free reference everywhere. In free
expansion (LC1) the duct lengthens and only the non-uniform temperature distribution creates stress. When both end faces are held axially
(LC2), the axial growth is prevented; the end supports then react an axial compressive force $N$ equal to the integrated restrained thermal
strain times the axial stiffness of the section. This force is the column load that governs the buckling analysis of Chapter~\ref{ch:buckling}.

\section{Baseline Operating Condition}
The baseline operating point P00 is summarised in Table~\ref{tab:operating}. The inlet Reynolds number of about 30,000 places the flow well
inside the turbulent regime; the bulk temperature rise of about 69~K reduces the Reynolds number to 25,545 at the outlet.

%% generated by gen_tables.py - RE-ANALYSIS 2026
\begin{table}[htbp]
\centering
\small
\caption{Baseline operating condition P00 (selected inputs)}
\label{tab:operating}
\begin{tabular}{>{\raggedright\arraybackslash}p{0.45\linewidth}rl>{\raggedright\arraybackslash}p{0.25\linewidth}}
\toprule
Quantity & Value & Unit & Status \\
\midrule
Fluid & air & -- & selected \\
Inlet temperature $T_{in}$ & 300 & K & selected \\
Inlet velocity $V_{in}$ (uniform) & 23.5 & m/s & selected \\
Inlet turbulence intensity & 4.411 & \% & assumed (A-020b) \\
Outlet static pressure (gauge) & 0 & Pa & selected \\
Operating pressure $p_{op}$ & 101,325 & Pa & selected \\
Outer-wall heat flux $q''$ (uniform) & 8000 & W/m$^2$ & selected \\
Equivalent bore flux $q''D_o/D_i$ & 16,000 & W/m$^2$ & calculated \\
End faces & adiabatic & -- & assumed (A-011) \\
Inlet Reynolds number (CFD) & 29,958 & -- & simulated \\
\bottomrule
\end{tabular}
\par\smallskip\parbox{\linewidth}{\footnotesize\raggedright Source: \texttt{MASTER\_PROJECT\_DATA.csv} rows M002--M006 and M034; \texttt{baseline\_parameters.json}.}
\end{table}

